% !TeX program = pdflatex
% !BIB program = biber
\documentclass[12pt,a4paper]{article}

\usepackage[a4paper,margin=1in]{geometry}
\usepackage{amsmath,amssymb,amsthm,mathtools,mathrsfs}
\usepackage{tikz-cd}
\usepackage{ytableau}
\usepackage{enumitem}
\usepackage[hidelinks]{hyperref}
\usepackage[style=alphabetic,backend=biber,sorting=nyt,maxbibnames=99]{biblatex}
\addbibresource{schur_weyl.bib}

\allowdisplaybreaks
\ytableausetup{boxsize=1.25em}

\newtheorem{theorem}{Theorem}[section]
\newtheorem{proposition}[theorem]{Proposition}
\newtheorem{lemma}[theorem]{Lemma}
\newtheorem{corollary}[theorem]{Corollary}

\theoremstyle{definition}
\newtheorem{definition}[theorem]{Definition}
\newtheorem{example}[theorem]{Example}
\newtheorem{exercise}[theorem]{Exercise}

\theoremstyle{remark}
\newtheorem{remark}[theorem]{Remark}

\newcommand{\C}{\mathbb{C}}
\newcommand{\Z}{\mathbb{Z}}
\newcommand{\GL}{\operatorname{GL}}
\newcommand{\End}{\operatorname{End}}
\newcommand{\Hom}{\operatorname{Hom}}
\newcommand{\Sym}{\operatorname{Sym}}
\newcommand{\sgn}{\operatorname{sgn}}
\newcommand{\Span}{\operatorname{span}}
\newcommand{\tr}{\operatorname{tr}}
\newcommand{\id}{\operatorname{id}}
\newcommand{\Schur}[1]{\mathbb{S}_{#1}}

\title{An Introduction to Schur--Weyl Duality}
\author{Thesis Draft}
\date{\today}

\begin{document}

\maketitle

\begin{abstract}
Schur--Weyl duality explains a remarkable compatibility between two
apparently different kinds of symmetry.  If \(V\) is a complex vector space
and \(V^{\otimes k}\) is its \(k\)-fold tensor power, then the general linear
group \(\GL(V)\) acts by changing every tensor factor in the same way, while
the symmetric group \(S_k\) acts by permuting tensor positions.  These two
actions commute, and the resulting decomposition of \(V^{\otimes k}\) is
simultaneously controlled by irreducible representations of \(S_k\) and by
polynomial representations of \(\GL(V)\).

This thesis gives an introduction to the classical
Schur--Weyl decomposition.  The route is deliberately concrete: we begin with
finite-dimensional representations, Young diagrams, Young symmetrizers, and
Specht modules; then we construct Schur functors from tensor powers; finally
we use the double-centralizer principle to prove the main theorem.  Several
low-degree examples are included to make the decomposition computable rather
than merely formal.  A final section sketches the connection with Schur
polynomials and highest weights.
\end{abstract}

\tableofcontents

\section{Introduction}

Tensor products are among the most useful constructions in linear algebra.
They let one encode multilinear maps, build higher-degree functions from
linear data, and turn a vector space \(V\) into a family of larger spaces
\[
V,\quad V^{\otimes 2},\quad V^{\otimes 3},\quad \ldots .
\]
When \(V\) is a complex vector space of dimension \(n\), the tensor power
\(V^{\otimes k}\) has dimension \(n^k\).  At first sight this space looks too
large to understand directly.  Schur--Weyl duality says that it is large in a
highly organized way.

There are two natural symmetries on \(V^{\otimes k}\).  First, every
invertible linear map \(g\in \GL(V)\) acts diagonally:
\[
g\cdot(v_1\otimes \cdots \otimes v_k)
  =(gv_1)\otimes \cdots \otimes (gv_k).
\]
Second, the symmetric group \(S_k\) acts by permuting the tensor positions.
In this draft we use a right action and the convention
\((\sigma\tau)(i)=\sigma(\tau(i))\):
\[
(v_1\otimes \cdots \otimes v_k)\cdot \sigma
  =v_{\sigma(1)}\otimes \cdots \otimes v_{\sigma(k)}.
\]
The central observation is that these two actions commute.  Applying a
linear change of coordinates to every tensor factor and then permuting the
factors gives the same result as doing the permutation first and the linear
change second.

The theorem proved in this thesis identifies the simultaneous decomposition
that follows from this commuting pair of actions:
\[
V^{\otimes k}
  \cong
  \bigoplus_{\lambda\vdash k,\ \ell(\lambda)\le n}
  S^\lambda\otimes \Schur{\lambda}(V).
\]
Here \(\lambda\vdash k\) means that \(\lambda\) is a partition of \(k\),
\(\ell(\lambda)\) is the number of nonzero parts of \(\lambda\),
\(S^\lambda\) is the Specht module of the symmetric group \(S_k\), and
\(\Schur{\lambda}(V)\) is the Schur functor attached to \(\lambda\).  Thus
the same partition \(\lambda\) labels an irreducible object on the
symmetric-group side and an irreducible polynomial representation on the
general-linear side.

The main goal is expository.  We will not attempt to reproduce every proof
from the representation theory of symmetric groups.  Instead, we prove the
basic facts needed to understand the theorem, state the standard
classification theorem for Specht modules, prove the commuting action
calculation in full, and then use a double-centralizer theorem as the proof
spine.  This follows the classical viewpoint of Schur and Weyl while keeping
the presentation close to modern introductory references such as
\cite{FultonHarris1991,Fulton1997,Etingof2011}.

\subsection*{Notation}

Throughout this thesis:
\begin{itemize}[leftmargin=2.2em]
  \item \(V\) denotes a finite-dimensional complex vector space;
  \item \(n=\dim V\);
  \item \(k\) denotes the tensor-power degree;
  \item \(S_k\) denotes the symmetric group on \(\{1,\dots,k\}\);
  \item \(\lambda\vdash k\) means that \(\lambda\) is a partition of \(k\);
  \item \(S^\lambda\) denotes the Specht module indexed by \(\lambda\);
  \item \(\Schur{\lambda}(V)\) denotes the Schur functor indexed by
  \(\lambda\).
\end{itemize}
All vector spaces and algebras are over \(\C\) unless otherwise stated.

\section{Background on Representations}

This section recalls the representation-theoretic language used later.  The
main point is that finite-group representations over \(\C\) split into
irreducible pieces.  Schur--Weyl duality is, in one sense, a very precise
description of how such irreducible pieces occur inside a tensor power.

\subsection{Representations and irreducibility}

\begin{definition}
Let \(G\) be a group.  A finite-dimensional representation of \(G\) is a
finite-dimensional complex vector space \(W\) together with a group
homomorphism
\[
\rho:G\longrightarrow \GL(W).
\]
We often write \(g\cdot w\) instead of \(\rho(g)w\).
\end{definition}

\begin{definition}
A subspace \(U\subseteq W\) is \(G\)-invariant if \(g\cdot u\in U\) for every
\(g\in G\) and every \(u\in U\).  A nonzero representation \(W\) is
irreducible if its only invariant subspaces are \(0\) and \(W\).
\end{definition}

\begin{example}
Let \(S_k\) act on \(\C^k\) by permuting the standard basis vectors
\(e_1,\dots,e_k\).  The line \(\C(e_1+\cdots+e_k)\) is invariant.  Its
complement
\[
\{(a_1,\dots,a_k)\in \C^k:a_1+\cdots+a_k=0\}
\]
is also invariant.  For \(k\ge 3\), this second representation is the
standard representation of \(S_k\).
\end{example}

\begin{definition}
If \(G\) acts on \(W\), then the group algebra \(\C[G]\) is the complex vector
space with basis \(\{e_g:g\in G\}\), with multiplication determined by
\(e_g e_h=e_{gh}\).  A representation of \(G\) is equivalently a left
\(\C[G]\)-module.
\end{definition}

The group algebra packages all group elements into one algebra.  This is
especially useful because the symmetric group action on \(V^{\otimes k}\)
extends linearly to an action of \(\C[S_k]\).

\subsection{Maschke's theorem and Schur's lemma}

\begin{theorem}[Maschke's theorem]
Let \(G\) be a finite group and let \(W\) be a finite-dimensional complex
representation of \(G\).  Then \(W\) is a direct sum of irreducible
representations.
\end{theorem}

\begin{proof}
It is enough to show that every invariant subspace has an invariant
complement.  Let \(U\subseteq W\) be invariant and choose any linear
projection \(p:W\to U\).  This projection need not commute with the action of
\(G\), so we average it:
\[
P(w)=\frac{1}{|G|}\sum_{g\in G}g\cdot p(g^{-1}\cdot w).
\]
Then \(P\) is still a projection onto \(U\), and the averaging makes it
\(G\)-equivariant: \(P(h\cdot w)=h\cdot P(w)\).  Therefore \(\ker P\) is
\(G\)-invariant and
\[
W=U\oplus \ker P.
\]
Repeating this argument inductively on dimension decomposes \(W\) into
irreducible summands.
\end{proof}

\begin{lemma}[Schur's lemma]
Let \(U\) and \(W\) be irreducible finite-dimensional representations of a
group \(G\).  If \(T:U\to W\) is \(G\)-equivariant, then either \(T=0\) or
\(T\) is an isomorphism.  In particular, if \(U=W\), then every
\(G\)-equivariant endomorphism of \(U\) is scalar.
\end{lemma}

\begin{proof}
The kernel of \(T\) is an invariant subspace of \(U\), and the image of \(T\)
is an invariant subspace of \(W\).  Since \(U\) and \(W\) are irreducible,
the only possibilities are \(\ker T=U\), giving \(T=0\), or
\(\ker T=0\) and \(\operatorname{im}T=W\), giving an isomorphism.

For the last statement, let \(T:U\to U\) be equivariant.  Since \(\C\) is
algebraically closed, \(T\) has an eigenvalue \(\alpha\).  Then
\(T-\alpha \id_U\) is equivariant and not injective, hence it is zero.
\end{proof}

\subsection{Characters and multiplicities}

\begin{definition}
The character of a finite-dimensional representation \(\rho:G\to\GL(W)\) is
the function \(\chi_W:G\to \C\) defined by
\[
\chi_W(g)=\tr(\rho(g)).
\]
\end{definition}

Characters are important because they remember multiplicities.  If
\[
W\cong \bigoplus_i m_i U_i
\]
is a decomposition into irreducible representations, then
\[
\chi_W=\sum_i m_i\chi_{U_i}.
\]
For finite groups over \(\C\), irreducible characters form an orthonormal
basis of the space of class functions.  We will not need the full machinery
of character theory, but the philosophy is helpful: representations can be
understood by decomposing them into irreducibles and tracking how many times
each one appears.

\begin{definition}
Let \(G\) act on \(W\).  The commutant, or centralizer algebra, of the action
is
\[
\End_G(W)=\{T\in \End(W):T(g\cdot w)=g\cdot T(w)
\text{ for all }g\in G,\ w\in W\}.
\]
\end{definition}

The commutant records the endomorphisms that are invisible to the group
action.  Schur--Weyl duality says that, for \(W=V^{\otimes k}\), the
endomorphisms invisible to the \(\GL(V)\)-action come from \(S_k\), and the
endomorphisms invisible to the \(S_k\)-action come from \(\GL(V)\).

\section{Partitions, Young Diagrams, and Specht Modules}

The irreducible representations of \(S_k\) are indexed by partitions of
\(k\).  This section introduces the combinatorial objects that make that
classification concrete.

\subsection{Partitions and Young diagrams}

\begin{definition}
A partition of \(k\), written \(\lambda\vdash k\), is a sequence
\[
\lambda=(\lambda_1,\lambda_2,\dots,\lambda_r)
\]
of positive integers such that
\[
\lambda_1\ge \lambda_2\ge \cdots\ge \lambda_r>0,
\qquad
\lambda_1+\cdots+\lambda_r=k.
\]
The number \(r\) is the length of \(\lambda\), denoted \(\ell(\lambda)\).
\end{definition}

We draw partitions using English Young diagrams: the first row has
\(\lambda_1\) boxes, the second row has \(\lambda_2\) boxes, and so on.  For
example, the partition \((3,1)\vdash 4\) has diagram
\[
\ydiagram{3,1}
\]
and the partition \((2,1)\vdash 3\) has diagram
\[
\ydiagram{2,1}.
\]

\begin{definition}
A Young tableau of shape \(\lambda\vdash k\) is a filling of the boxes of the
Young diagram of \(\lambda\) with the numbers \(1,\dots,k\), each used once.
The tableau is standard if its entries increase along each row from left to
right and down each column from top to bottom.
\end{definition}

\begin{example}
The tableau
\[
\begin{ytableau}
1 & 2\\
3
\end{ytableau}
\]
is a standard tableau of shape \((2,1)\).  The tableau
\[
\begin{ytableau}
1 & 3\\
2
\end{ytableau}
\]
is also standard.  These two standard tableaux already hint at the fact that
the Specht module \(S^{(2,1)}\) has dimension \(2\).
\end{example}

\subsection{Young symmetrizers}

Fix a Young tableau \(T\) of shape \(\lambda\vdash k\).  Two subgroups of
\(S_k\) are attached to \(T\):
\begin{itemize}[leftmargin=2.2em]
  \item the row group \(R_T\), consisting of permutations that preserve each
  row of \(T\) as a set;
  \item the column group \(C_T\), consisting of permutations that preserve
  each column of \(T\) as a set.
\end{itemize}

\begin{definition}
Inside the group algebra \(\C[S_k]\), define
\[
a_T=\sum_{\sigma\in R_T}\sigma,
\qquad
b_T=\sum_{\tau\in C_T}\sgn(\tau)\tau,
\qquad
c_T=a_Tb_T.
\]
The element \(c_T\) is called a Young symmetrizer.
\end{definition}

The element \(a_T\) symmetrizes along rows, while \(b_T\) antisymmetrizes
along columns.  The order matters.  In many references one also considers
\(b_Ta_T\), and the resulting modules are naturally isomorphic after the
usual left-right convention is handled carefully.

\begin{example}
For the tableau
\[
T=
\begin{ytableau}
1 & 2\\
3
\end{ytableau},
\]
the row group is \(R_T=\{1,(12)\}\), while the column group is
\(C_T=\{1,(13)\}\).  Hence
\[
a_T=1+(12),\qquad b_T=1-(13),
\]
and
\[
c_T=(1+(12))(1-(13))
  =1-(13)+(12)-(132).
\]
This element first alternates in the first column and then symmetrizes in the
first row.
\end{example}

\subsection{Specht modules}

\begin{definition}
Let \(\lambda\vdash k\), and choose a tableau \(T\) of shape \(\lambda\).  The
Specht module \(S^\lambda\) is the irreducible \(S_k\)-module constructed
from the Young symmetrizer \(c_T\).  Equivalently, it can be realized as the
left module generated by a Young symmetrizer, after the standard convention
for left and right actions is fixed.
\end{definition}

The construction above depends on choosing a tableau, but the resulting
irreducible representation depends only on the shape \(\lambda\), up to
isomorphism.  This is one of the central results in the representation theory
of the symmetric group.

\begin{theorem}[Specht module classification]
The irreducible complex representations of \(S_k\) are indexed by partitions
\(\lambda\vdash k\).  The representation corresponding to \(\lambda\) is the
Specht module \(S^\lambda\).
\end{theorem}

We use this theorem as a standard input.  A full proof is beautiful but long:
it requires comparing Young symmetrizers, dominance order on partitions, and
the structure of the group algebra \(\C[S_k]\).  For the purposes of
Schur--Weyl duality, the important consequence is that the \(S_k\)-side of
the decomposition of \(V^{\otimes k}\) is completely indexed by partitions of
\(k\).

\begin{theorem}[Hook-length formula]
For a box \(x\) in the Young diagram of \(\lambda\), let \(h(x)\) be its hook
length: the number of boxes weakly to the right of \(x\) in the same row plus
the number of boxes strictly below \(x\) in the same column.  Then
\[
\dim S^\lambda
  =\frac{k!}{\prod_{x\in\lambda}h(x)}.
\]
\end{theorem}

\begin{example}
For \(\lambda=(2,1)\), the hook lengths are
\[
\begin{ytableau}
3 & 1\\
1
\end{ytableau}
\]
so
\[
\dim S^{(2,1)}=\frac{3!}{3\cdot 1\cdot 1}=2.
\]
The other partitions of \(3\) are \((3)\) and \((1,1,1)\).  Their Specht
modules are the trivial and sign representations, both one-dimensional.
\end{example}

\begin{exercise}
Use the hook-length formula to compute the dimensions of the Specht modules
for the five partitions of \(4\):
\[
(4),\quad (3,1),\quad (2,2),\quad (2,1,1),\quad (1,1,1,1).
\]
\end{exercise}

\section{Tensor Powers and Schur Functors}

We now return to tensor powers.  This section defines the two commuting
actions on \(V^{\otimes k}\), proves that they commute, and introduces Schur
functors as the multiplicity spaces attached to Specht modules.

\subsection{The two actions}

Let \(V\) be a complex vector space of dimension \(n\).  The group \(\GL(V)\)
acts on \(V^{\otimes k}\) by
\[
g\cdot(v_1\otimes\cdots\otimes v_k)
  =(gv_1)\otimes\cdots\otimes(gv_k).
\]
The symmetric group \(S_k\) acts on the right by
\[
(v_1\otimes\cdots\otimes v_k)\cdot\sigma
  =v_{\sigma(1)}\otimes\cdots\otimes v_{\sigma(k)}.
\]
This action extends linearly to an action of the group algebra \(\C[S_k]\).

\begin{proposition}
The left action of \(\GL(V)\) and the right action of \(S_k\) on
\(V^{\otimes k}\) commute.
\end{proposition}

\begin{proof}
It is enough to check pure tensors.  Let \(g\in \GL(V)\), let
\(\sigma\in S_k\), and let \(v_1,\dots,v_k\in V\).  Then
\[
\begin{aligned}
\bigl(g\cdot(v_1\otimes\cdots\otimes v_k)\bigr)\cdot\sigma
&=((gv_1)\otimes\cdots\otimes(gv_k))\cdot\sigma\\
&=(gv_{\sigma(1)})\otimes\cdots\otimes(gv_{\sigma(k)}).
\end{aligned}
\]
On the other hand,
\[
\begin{aligned}
g\cdot\bigl((v_1\otimes\cdots\otimes v_k)\cdot\sigma\bigr)
&=g\cdot(v_{\sigma(1)}\otimes\cdots\otimes v_{\sigma(k)})\\
&=(gv_{\sigma(1)})\otimes\cdots\otimes(gv_{\sigma(k)}).
\end{aligned}
\]
The two expressions agree, so the actions commute.
\end{proof}

The situation can be summarized by the following commuting square:
\[
\begin{tikzcd}[column sep=huge,row sep=huge]
V^{\otimes k}
  \arrow[r, "{\text{right action of }S_k}"]
  \arrow[d, "{\text{left action of }\GL(V)}"']
&
V^{\otimes k}
  \arrow[d, "{\text{left action of }\GL(V)}"]
\\
V^{\otimes k}
  \arrow[r, "{\text{right action of }S_k}"']
&
V^{\otimes k}.
\end{tikzcd}
\]

\subsection{Definition of the Schur functor}

Because \(V^{\otimes k}\) is a right \(\C[S_k]\)-module and \(S^\lambda\) is
a left \(\C[S_k]\)-module, we may form their tensor product over
\(\C[S_k]\).

\begin{definition}
For \(\lambda\vdash k\), define the Schur functor
\[
\Schur{\lambda}(V)
  :=V^{\otimes k}\otimes_{\C[S_k]}S^\lambda.
\]
\end{definition}

There are several equivalent descriptions.  One can also view
\(\Schur{\lambda}(V)\) as a multiplicity space:
\[
\Schur{\lambda}(V)\cong \Hom_{S_k}(S^\lambda,V^{\otimes k}),
\]
with the natural \(\GL(V)\)-action.  Alternatively, after choosing a Young
symmetrizer \(c_T\), one can identify \(\Schur{\lambda}(V)\) with the image
of \(c_T\) acting on \(V^{\otimes k}\), up to harmless normalizing constants
and left-right conventions.

\begin{example}
The partition \((k)\) has one row.  The corresponding Young symmetrizer
symmetrizes all tensor positions, and
\[
\Schur{(k)}(V)\cong \Sym^k(V).
\]
The partition \((1^k)\) has one column.  The corresponding Young symmetrizer
antisymmetrizes all tensor positions, and
\[
\Schur{(1^k)}(V)\cong \bigwedge^k V.
\]
\end{example}

\subsection{Basic properties}

\begin{proposition}
There are natural isomorphisms
\[
\Schur{(k)}(V)\cong \Sym^k(V),
\qquad
\Schur{(1^k)}(V)\cong \bigwedge^k V.
\]
\end{proposition}

\begin{proof}
For \((k)\), the row group is all of \(S_k\) and the column group is trivial.
The Young symmetrizer is therefore the full symmetrizer
\[
\sum_{\sigma\in S_k}\sigma.
\]
Its image consists exactly of symmetric tensors, which represent
\(\Sym^k(V)\).

For \((1^k)\), the row group is trivial and the column group is all of
\(S_k\).  The Young symmetrizer is the alternating operator
\[
\sum_{\sigma\in S_k}\sgn(\sigma)\sigma.
\]
Its image consists exactly of alternating tensors, which represent
\(\bigwedge^k V\).
\end{proof}

\begin{proposition}
If \(\ell(\lambda)>n=\dim V\), then \(\Schur{\lambda}(V)=0\).
\end{proposition}

\begin{proof}
If \(\ell(\lambda)>n\), then the Young diagram of \(\lambda\) has a column
with more than \(n\) boxes.  The Young symmetrizer defining
\(\Schur{\lambda}(V)\) includes an antisymmetrization along this column.  But
there is no nonzero alternating tensor of degree \(n+1\) or higher in an
\(n\)-dimensional vector space.  Thus the corresponding image in
\(V^{\otimes k}\) is zero.
\end{proof}

\begin{remark}
The length restriction \(\ell(\lambda)\le n\) is one of the main differences
between the two sides of Schur--Weyl duality.  Every partition of \(k\) gives
a Specht module for \(S_k\), but only partitions with at most \(n\) rows give
nonzero Schur functors on an \(n\)-dimensional vector space.
\end{remark}

\section{Schur--Weyl Duality}

The proof of Schur--Weyl duality is best understood through centralizers.
The two actions on \(V^{\otimes k}\) commute, so each action lands inside the
centralizer of the other.  Schur--Weyl duality says that these inclusions are
as large as possible.

\subsection{The double-centralizer principle}

\begin{theorem}[Double-centralizer theorem]
Let \(A\) be a finite-dimensional semisimple algebra over \(\C\), and let
\(W\) be a finite-dimensional \(A\)-module.  Suppose
\[
W\cong \bigoplus_i U_i\otimes M_i
\]
as an \(A\)-module, where the \(U_i\) are pairwise nonisomorphic irreducible
\(A\)-modules and \(A\) acts on \(U_i\otimes M_i\) only through \(U_i\).
Then
\[
\End_A(W)\cong \bigoplus_i \End(M_i).
\]
Moreover, the centralizer of \(\End_A(W)\) inside \(\End(W)\) is
\[
\bigoplus_i \End(U_i),
\]
acting on the \(U_i\)-factors.
\end{theorem}

\begin{proof}
By Schur's lemma, an \(A\)-equivariant map from \(U_i\otimes M_i\) to
\(U_j\otimes M_j\) must be zero when \(i\ne j\).  When \(i=j\), the map must
be the identity on \(U_i\), up to scalar on the irreducible factor, tensored
with an arbitrary linear map on \(M_i\).  This gives
\[
\End_A(W)\cong \bigoplus_i \End(M_i).
\]
Taking the centralizer once more exchanges the roles of the irreducible
factors \(U_i\) and the multiplicity spaces \(M_i\), giving the second
statement.
\end{proof}

This theorem is the algebraic reason that commuting actions produce a clean
decomposition.  One action decomposes the space into irreducibles; the other
action is forced to operate on the multiplicity spaces.

\subsection{The centralizer statement}

Let
\[
\rho:\GL(V)\to \GL(V^{\otimes k})
\]
be the diagonal action, and let
\[
\pi:\C[S_k]\to \End(V^{\otimes k})
\]
be the linear extension of the place-permutation action.  Since the two
actions commute, we have inclusions
\[
\rho(\GL(V))\subseteq \End_{\C[S_k]}(V^{\otimes k}),
\qquad
\pi(\C[S_k])\subseteq \End_{\GL(V)}(V^{\otimes k}).
\]
The classical centralizer theorem says that these inclusions become
equalities if the left side is interpreted linearly:
\[
\Span_\C(\rho(\GL(V)))=\End_{\C[S_k]}(V^{\otimes k}),
\qquad
\pi(\C[S_k])=\End_{\GL(V)}(V^{\otimes k}).
\]

We record this as a standard theorem.  A full proof can be given using
matrix coefficient calculations, polynomial density, or the representation
theory of \(\mathfrak{gl}(V)\).

\begin{theorem}[Schur--Weyl centralizer theorem]
On \(V^{\otimes k}\), the image of \(\C[S_k]\) and the linear span of the
image of \(\GL(V)\) are mutual centralizers in \(\End(V^{\otimes k})\).
\end{theorem}

\subsection{The main theorem}

\begin{theorem}[Schur--Weyl duality]
Let \(V\) be a complex vector space of dimension \(n\).  For every integer
\(k\ge 0\), there is an isomorphism of \(S_k\times \GL(V)\)-modules
\[
V^{\otimes k}
  \cong
  \bigoplus_{\lambda\vdash k,\ \ell(\lambda)\le n}
  S^\lambda\otimes \Schur{\lambda}(V).
\]
Furthermore, for each partition \(\lambda\vdash k\) with
\(\ell(\lambda)\le n\), the \(\GL(V)\)-module \(\Schur{\lambda}(V)\) is
irreducible.
\end{theorem}

\begin{proof}
By Maschke's theorem, \(V^{\otimes k}\), viewed as a representation of the
finite group \(S_k\), decomposes as a direct sum of irreducible
\(S_k\)-modules:
\[
V^{\otimes k}\cong \bigoplus_{\lambda\vdash k}S^\lambda\otimes M_\lambda,
\]
where
\[
M_\lambda\cong \Hom_{S_k}(S^\lambda,V^{\otimes k})
\]
is the multiplicity space.  The Specht classification tells us that the
irreducible \(S_k\)-modules are exactly the \(S^\lambda\).

Because the \(\GL(V)\)-action commutes with the \(S_k\)-action, \(\GL(V)\)
acts on each multiplicity space \(M_\lambda\).  By definition, this
multiplicity space is the Schur functor \(\Schur{\lambda}(V)\).  Thus
\[
V^{\otimes k}\cong
\bigoplus_{\lambda\vdash k}S^\lambda\otimes \Schur{\lambda}(V)
\]
as a module for the commuting pair of actions.

If \(\ell(\lambda)>n\), then \(\Schur{\lambda}(V)=0\), as proved earlier.
Hence only partitions with at most \(n\) parts remain.  Finally, the
double-centralizer theorem and the Schur--Weyl centralizer theorem imply that
the nonzero multiplicity spaces for the \(S_k\)-decomposition are
irreducible modules for the commuting \(\GL(V)\)-action.  Therefore
\(\Schur{\lambda}(V)\) is irreducible whenever \(\ell(\lambda)\le n\).
\end{proof}

\begin{corollary}
As a representation of \(\GL(V)\) alone, the multiplicity of
\(\Schur{\lambda}(V)\) in \(V^{\otimes k}\) is \(\dim S^\lambda\).
\end{corollary}

\begin{proof}
If we forget the \(S_k\)-action in
\[
V^{\otimes k}
  \cong
  \bigoplus_{\lambda\vdash k,\ \ell(\lambda)\le n}
  S^\lambda\otimes \Schur{\lambda}(V),
\]
then the factor \(S^\lambda\) becomes only a multiplicity space of dimension
\(\dim S^\lambda\).
\end{proof}

\section{Examples and Computations}

The examples in this section illustrate how the theorem decomposes familiar
tensor powers.  They also show why the dimensions on both sides match.

\subsection{Degree two}

For \(k=2\), the partitions are \((2)\) and \((1,1)\).  The corresponding
Schur functors are
\[
\Schur{(2)}(V)=\Sym^2(V),
\qquad
\Schur{(1,1)}(V)=\bigwedge^2V.
\]
The two Specht modules are one-dimensional: \(S^{(2)}\) is the trivial
representation of \(S_2\), and \(S^{(1,1)}\) is the sign representation.
Therefore Schur--Weyl duality gives
\[
V^{\otimes 2}\cong \Sym^2(V)\oplus \bigwedge^2V
\]
as a \(\GL(V)\)-module.

This decomposition can be proved directly.  Define
\[
p_+(v\otimes w)=\frac{1}{2}(v\otimes w+w\otimes v),
\qquad
p_-(v\otimes w)=\frac{1}{2}(v\otimes w-w\otimes v).
\]
Then \(p_+\) and \(p_-\) are projections,
\[
p_+^2=p_+,\qquad p_-^2=p_-,
\]
and
\[
p_++p_-=\id_{V^{\otimes 2}}.
\]
The image of \(p_+\) is \(\Sym^2(V)\), and the image of \(p_-\) is
\(\bigwedge^2V\).  Hence
\[
V^{\otimes 2}=\operatorname{im}p_+\oplus \operatorname{im}p_-.
\]

\subsection{Degree three}

For \(k=3\), the partitions are
\[
(3),\qquad (2,1),\qquad (1,1,1).
\]
The Specht dimensions are
\[
\dim S^{(3)}=1,\qquad
\dim S^{(2,1)}=2,\qquad
\dim S^{(1,1,1)}=1.
\]
If \(n=\dim V\ge 3\), then all three Schur functors are nonzero, and
Schur--Weyl duality gives, as a \(\GL(V)\)-module,
\[
V^{\otimes 3}
  \cong
  \Sym^3(V)\oplus
  \Schur{(2,1)}(V)^{\oplus 2}
  \oplus
  \bigwedge^3V.
\]
If \(n=2\), the last term vanishes because \(\bigwedge^3V=0\).

\begin{example}
Let \(V=\C^3\).  Then
\[
\dim \Sym^3(\C^3)=\binom{3+3-1}{3}=\binom{5}{3}=10,
\]
and
\[
\dim \bigwedge^3\C^3=1.
\]
Using the hook-content formula stated below,
\[
\dim \Schur{(2,1)}(\C^3)=8.
\]
Therefore
\[
10+2\cdot 8+1=27=\dim(\C^3)^{\otimes 3}.
\]
\end{example}

\subsection{Dimension formulas}

The Specht side has the hook-length formula.  The general-linear side has a
closely related formula.

\begin{theorem}[Hook-content formula]
Let \(\lambda\vdash k\) and \(V=\C^n\).  If \(\ell(\lambda)\le n\), then
\[
\dim \Schur{\lambda}(\C^n)
  =
  \prod_{(i,j)\in \lambda}
  \frac{n+j-i}{h(i,j)},
\]
where \(h(i,j)\) is the hook length of the box in row \(i\), column \(j\).
\end{theorem}

\begin{example}
For \(\lambda=(2,1)\), the boxes are \((1,1)\), \((1,2)\), and \((2,1)\).
Their hook lengths are \(3,1,1\).  Hence
\[
\dim \Schur{(2,1)}(\C^n)
  =
  \frac{n}{3}(n+1)(n-1)
  =
  \frac{n(n^2-1)}{3}.
\]
For \(n=2\), this dimension is \(2\); for \(n=3\), it is \(8\).
\end{example}

\begin{example}
Take \(V=\C^2\) and \(k=4\).  Only partitions of \(4\) with at most two rows
survive:
\[
(4),\qquad (3,1),\qquad (2,2).
\]
The hook-length formula gives
\[
\dim S^{(4)}=1,\qquad
\dim S^{(3,1)}=3,\qquad
\dim S^{(2,2)}=2.
\]
The hook-content formula gives
\[
\dim \Schur{(4)}(\C^2)=5,\qquad
\dim \Schur{(3,1)}(\C^2)=3,\qquad
\dim \Schur{(2,2)}(\C^2)=1.
\]
Thus
\[
1\cdot 5+3\cdot 3+2\cdot 1=16=\dim(\C^2)^{\otimes 4}.
\]
\end{example}

\subsection{A worked Young symmetrizer example}

Let \(\lambda=(2,1)\) and choose
\[
T=
\begin{ytableau}
1 & 2\\
3
\end{ytableau}.
\]
The Young symmetrizer is
\[
c_T=1-(13)+(12)-(132).
\]
Applying this to a pure tensor \(u\otimes v\otimes w\) gives
\[
\begin{aligned}
(u\otimes v\otimes w)c_T
&=u\otimes v\otimes w
  -w\otimes v\otimes u
  +v\otimes u\otimes w
  -w\otimes u\otimes v.
\end{aligned}
\]
The expression is symmetric in the first-row positions \(1,2\) and
alternating in the first-column positions \(1,3\), according to the tableau.
The span of such tensors is a concrete model for
\(\Schur{(2,1)}(V)\), up to the conventional choice of left or right
symmetrizer.

\section{Characters and Highest Weights}

The previous sections prove the structural decomposition.  This final
section briefly explains how the same representations are described by
characters and highest weights.  The point is not to develop the full theory
of Lie algebras, but to show where the notation \(\Schur{\lambda}(V)\) comes
from and why partitions are natural on the \(\GL(V)\)-side as well.

\subsection{Polynomial characters}

Let \(V=\C^n\), and let \(g\in \GL_n(\C)\) be diagonal with eigenvalues
\[
x_1,\dots,x_n.
\]
Then \(g\) acts diagonally on \(V^{\otimes k}\), and its trace is
\[
(x_1+\cdots+x_n)^k.
\]
Schur--Weyl duality decomposes this character into irreducible polynomial
characters:
\[
(x_1+\cdots+x_n)^k
  =
  \sum_{\lambda\vdash k,\ \ell(\lambda)\le n}
  (\dim S^\lambda)s_\lambda(x_1,\dots,x_n),
\]
where \(s_\lambda\) is the Schur polynomial indexed by \(\lambda\).

\begin{definition}
The Schur polynomial \(s_\lambda(x_1,\dots,x_n)\) is the character of the
\(\GL_n(\C)\)-representation \(\Schur{\lambda}(\C^n)\):
\[
\chi_{\Schur{\lambda}(\C^n)}(g)=s_\lambda(x_1,\dots,x_n)
\]
when \(g\) has eigenvalues \(x_1,\dots,x_n\).
\end{definition}

This definition is representation-theoretic.  There are also purely
combinatorial definitions of Schur polynomials, for example as generating
functions over semistandard Young tableaux.  The equality between these
viewpoints is one of the main bridges between representation theory and
algebraic combinatorics.

\subsection{Highest weights}

Let \(T\subseteq \GL_n(\C)\) be the subgroup of diagonal matrices.  A vector
in a representation \(W\) has weight \(\mu=(\mu_1,\dots,\mu_n)\in\Z^n\) if
\[
\operatorname{diag}(t_1,\dots,t_n)\cdot w
  =
  t_1^{\mu_1}\cdots t_n^{\mu_n}w.
\]
Polynomial irreducible representations of \(\GL_n(\C)\) are classified by
dominant weights
\[
\lambda_1\ge \lambda_2\ge \cdots\ge \lambda_n\ge 0.
\]
These are exactly partitions with at most \(n\) parts.

\begin{theorem}
The representation \(\Schur{\lambda}(\C^n)\) has highest weight
\(\lambda=(\lambda_1,\dots,\lambda_n)\), where trailing zeroes are added if
necessary.
\end{theorem}

This theorem places Schur--Weyl duality into the broader representation
theory of reductive groups and Lie algebras.  From that perspective, the
condition \(\ell(\lambda)\le n\) is simply the condition that \(\lambda\) be a
dominant polynomial weight for \(\GL_n(\C)\).

\begin{example}
For \(V=\C^3\), the representation \(\Schur{(2,1)}(V)\) has highest weight
\((2,1,0)\).  Its dimension is \(8\), and it appears twice in
\(V^{\otimes 3}\) because the corresponding Specht module \(S^{(2,1)}\) has
dimension \(2\).
\end{example}

\section{Conclusion}

Schur--Weyl duality describes \(V^{\otimes k}\) as a meeting point between
two representation theories.  The symmetric group \(S_k\) controls how
tensor positions may be permuted, while the general linear group \(\GL(V)\)
controls how the vector space \(V\) itself may be changed.  The actions
commute, and the double-centralizer theorem explains why this produces a
simultaneous decomposition.

The final decomposition
\[
V^{\otimes k}
  \cong
  \bigoplus_{\lambda\vdash k,\ \ell(\lambda)\le n}
  S^\lambda\otimes \Schur{\lambda}(V)
\]
summarizes the whole story.  The same partition \(\lambda\) indexes a Specht
module for \(S_k\) and a Schur functor for \(\GL(V)\).  The Specht dimension
measures the multiplicity of the corresponding \(\GL(V)\)-module, while the
dimension of \(V\) determines which partitions survive.

Several natural directions remain beyond this introductory draft.  One may
study the Robinson--Schensted--Knuth correspondence, develop the full theory
of symmetric functions, prove the complete highest-weight classification, or
replace \(\GL(V)\) by orthogonal and symplectic groups, where Brauer algebras
take the role played here by symmetric groups.  These extensions all preserve
the same guiding idea: tensor powers become understandable when their
commuting symmetries are studied together.

\printbibliography

\end{document}
